\section[Results]{Measured Peak Heat Release Rates}
\label{sec:results}

The ceiling jet under a $4\units{m}$ ceiling was calculated with Alpert's correlations~\cite{alpert1972}, using the
fast design fire of Section~\ref{sec:design-fires}. Figure~\ref{fig:ceilingjet} shows the temperature $3\units{m}$
from the fire's axis. A sprinkler rated at $68\units{^\circ C}$ with a response time index of
$100\units{(m\,s)^{1/2}}$ activates after about 184 seconds, when the fire releases about $1.6\units{MW}$.

\begin{figure}[htb]
  \centering
  \includegraphics[width=0.7\textwidth]{ceiling_jet}
  \caption{Ceiling-jet temperature $3\units{m}$ from the axis of a fast t-squared fire under a $4\units{m}$ ceiling.}
  \label{fig:ceilingjet}
\end{figure}

The measured sofa fire reached its peak sooner than any design fire: as Figure 3 shows, the peak heat release rate
of the sofa was about 2.9 MW. Detailed models of the same room, run with the Fire Dynamics
Simulator~\cite{mcgrattan2023fds, nist2026}, agree with the design fire to within 10\% until the sprinkler activates.
Their results depend on the grid size, as the technical reference guide notes~\cite{fdsguide2023}, and on the soot
yield that is chosen for the heat release rate in
in a compartment of this size.
